\documentclass[11pt,a4paper]{article}
\usepackage[utf8]{inputenc}
\usepackage[T1]{fontenc}
\usepackage{amsmath,amssymb,amsthm}
\usepackage{physics}
\usepackage{geometry}
\usepackage{hyperref}
\usepackage{xcolor}
\usepackage{tcolorbox}
\usepackage{bm}

\geometry{margin=2.5cm}

\hypersetup{
    colorlinks=true,
    linkcolor=blue,
    citecolor=blue,
    urlcolor=blue
}

\newcommand{\cc}{\mathrm{c.c.}}
\newcommand{\re}{\mathrm{Re}}
\newcommand{\im}{\mathrm{Im}}
\newcommand{\R}{\mathbb{R}}
\newcommand{\C}{\mathbb{C}}
\newcommand{\half}{\tfrac{1}{2}}

\theoremstyle{definition}
\newtheorem{definition}{Definition}
\newtheorem{remark}{Remark}

\title{\textbf{Analytic Matrix Elements for Gaussian Basis Functions}\\[0.5em]
\large Complete Derivation of Overlap, Kinetic, and Potential Integrals\\for the Rothe Time-Stepping Method}
\author{Generated Documentation}
\date{\today}

\begin{document}

\maketitle
\tableofcontents
\newpage

%%%%%%%%%%%%%%%%%%%%%%%%%%%%%%%%%%%%%%%%%%%%%%%%%%%%%%%%%%%%%%%%%%%%%%%%%%%%%%%
\section{Introduction and Notation}
%%%%%%%%%%%%%%%%%%%%%%%%%%%%%%%%%%%%%%%%%%%%%%%%%%%%%%%%%%%%%%%%%%%%%%%%%%%%%%%

This document provides a complete derivation of all matrix elements used in the JAX-based Rothe solver for the time-dependent Schrödinger equation. We work in $D$ spatial dimensions with $n$ non-orthogonal Gaussian basis functions.

\subsection{Gaussian Basis Functions}

\begin{definition}[Generalized Gaussian]
Each basis function $g_i(\bm{x})$ for $i = 1, \ldots, n$ is a product of one-dimensional Gaussians:
\begin{equation}
\boxed{
g_i(\bm{x}) = \prod_{d=1}^{D} \phi_i^{(d)}(x_d), \quad
\phi_i^{(d)}(x) = \exp\left[ -\alpha_i^{(d)} (x - \mu_i^{(d)})^2 + i p_i^{(d)} x \right]
}
\end{equation}
where the complex width parameter is
\begin{equation}
\alpha_i^{(d)} = (a_i^{(d)})^2 + i b_i^{(d)}, \quad a_i^{(d)} > 0.
\end{equation}
\end{definition}

The parameters stored for each Gaussian are:
\begin{itemize}
    \item $a_i^{(d)} > 0$: width parameter (real, stored as $a$ so that $\re(\alpha) = a^2 > 0$)
    \item $b_i^{(d)} \in \R$: quadratic phase (chirp)
    \item $\mu_i^{(d)} \in \R$: center position
    \item $p_i^{(d)} \in \R$: momentum (linear phase)
\end{itemize}

These are stored in an array of shape \texttt{(n, D, 4)} with order \texttt{(a, b, mu, p)}.

\subsection{Wavefunction Ansatz}

The wavefunction is a linear combination:
\begin{equation}
\Psi(\bm{x}, t) = \sum_{i=1}^{n} c_i(t) \, g_i(\bm{x}; \bm{\theta}_i(t)),
\end{equation}
where $c_i \in \C$ are complex coefficients and $\bm{\theta}_i = (a_i, b_i, \mu_i, p_i)$ are the Gaussian parameters.

%%%%%%%%%%%%%%%%%%%%%%%%%%%%%%%%%%%%%%%%%%%%%%%%%%%%%%%%%%%%%%%%%%%%%%%%%%%%%%%
\section{Overlap Matrix $S$}
%%%%%%%%%%%%%%%%%%%%%%%%%%%%%%%%%%%%%%%%%%%%%%%%%%%%%%%%%%%%%%%%%%%%%%%%%%%%%%%

\subsection{One-Dimensional Gaussian Overlap}

Consider the overlap of two 1D Gaussians:
\begin{equation}
\langle \phi_i | \phi_j \rangle = \int_{-\infty}^{\infty} \phi_i^*(x) \phi_j(x) \, dx.
\end{equation}

Expanding the integrand:
\begin{align}
\phi_i^*(x) \phi_j(x) &= \exp\left[ -\alpha_i^* (x - \mu_i)^2 - i p_i x \right] \cdot \exp\left[ -\alpha_j (x - \mu_j)^2 + i p_j x \right]
\end{align}

Expanding the quadratics:
\begin{align}
-\alpha_i^*(x-\mu_i)^2 - \alpha_j(x-\mu_j)^2 &= -(\alpha_i^* + \alpha_j)x^2 + 2(\alpha_i^*\mu_i + \alpha_j\mu_j)x - \alpha_i^*\mu_i^2 - \alpha_j\mu_j^2
\end{align}

Adding the linear phase terms $-ip_i x + ip_j x = i(p_j - p_i)x$:
\begin{align}
\phi_i^*(x)\phi_j(x) &= \exp\Big[ -\underbrace{(\alpha_i^* + \alpha_j)}_{\Gamma} x^2 + \underbrace{2(\alpha_i^*\mu_i + \alpha_j\mu_j) + i(p_j - p_i)}_{\text{linear coefficient}} \, x \\
&\qquad\quad + \underbrace{(-\alpha_i^*\mu_i^2 - \alpha_j\mu_j^2)}_{\text{constant}} \Big] \nonumber
\end{align}

Define $A = \Gamma$, $B = 2(\alpha_i^*\mu_i + \alpha_j\mu_j) + i(p_j - p_i)$, $C = -\alpha_i^*\mu_i^2 - \alpha_j\mu_j^2$, so:
\begin{equation}
\phi_i^*(x)\phi_j(x) = \exp\left[ -Ax^2 + Bx + C \right]
\end{equation}

Completing the square: $-Ax^2 + Bx = -A\left(x - \frac{B}{2A}\right)^2 + \frac{B^2}{4A}$, giving:
\begin{equation}
\int_{-\infty}^{\infty} \phi_i^*(x)\phi_j(x)\,dx = \exp\left[C + \frac{B^2}{4A}\right] \cdot \sqrt{\frac{\pi}{A}} = \exp\left[C + \frac{B^2}{4A} - \half\ln A\right] \cdot \sqrt{\pi}
\end{equation}

The $\sqrt{\pi}$ factors are absorbed into normalization.

\begin{definition}[Pair Intermediates]
For a pair of Gaussians $(i, j)$ in dimension $d$, define:
\begin{align}
\Gamma_{ij}^{(d)} &= \alpha_i^{(d)*} + \alpha_j^{(d)} \label{eq:Gamma} \\[0.3em]
\tilde{\mu}_i^{(d)} &= \mu_i^{(d)} - \frac{i p_i^{(d)}}{2 \alpha_i^{(d)*}} \\[0.3em]
\tilde{\mu}_j^{(d)} &= \mu_j^{(d)} + \frac{i p_j^{(d)}}{2 \alpha_j^{(d)}} \\[0.3em]
\tilde{y}_{ij}^{(d)} &= \tilde{\mu}_i^{(d)} - \tilde{\mu}_j^{(d)} \\[0.3em]
\tilde{k}_{ij}^{(d)} &= \frac{\alpha_i^{(d)*} \alpha_j^{(d)}}{\Gamma_{ij}^{(d)}}
\end{align}
\end{definition}

The exponent contribution from dimension $d$ (excluding prefactors) is:
\begin{equation}
E_{ij}^{(d)} = -\frac{(p_i^{(d)})^2}{4\alpha_i^{(d)*}} - \frac{(p_j^{(d)})^2}{4\alpha_j^{(d)}} - \tilde{k}_{ij}^{(d)} (\tilde{y}_{ij}^{(d)})^2.
\end{equation}

\subsection{Full Overlap Matrix}

\begin{tcolorbox}[colback=blue!5,colframe=blue!50!black,title=Overlap Matrix $S$]
\begin{equation}
S_{ij} = \langle g_i | g_j \rangle = \exp\left[ \sum_{d=1}^{D} \left( E_{ij}^{(d)} - \half \ln \Gamma_{ij}^{(d)} \right) - \half \sum_{d=1}^{D} \ln(2 a_i^{(d)2}) - \half \sum_{d=1}^{D} \ln(2 a_j^{(d)2}) \right]
\end{equation}
\end{tcolorbox}

The last two terms are the diagonal normalization factors ensuring $S_{ii} = 1$ for all parameter values.

\textbf{Implementation note:} In the code, the overlap is computed as:
\begin{verbatim}
expo_core = sum_d(exponent_contrib - tilde_y^2 * tilde_k)
expo_pref = -0.5 * sum_d(log(Gamma))
diag_expo = -0.5 * sum_d(log(2 * a^2))
S = exp(expo_core + expo_pref - 0.5*(diag_expo_i + diag_expo_j))
\end{verbatim}

%%%%%%%%%%%%%%%%%%%%%%%%%%%%%%%%%%%%%%%%%%%%%%%%%%%%%%%%%%%%%%%%%%%%%%%%%%%%%%%
\section{Kinetic Energy Matrix $T$}
%%%%%%%%%%%%%%%%%%%%%%%%%%%%%%%%%%%%%%%%%%%%%%%%%%%%%%%%%%%%%%%%%%%%%%%%%%%%%%%

The kinetic energy operator is:
\begin{equation}
\hat{T} = -\frac{1}{2} \sum_{d=1}^{D} \frac{\partial^2}{\partial x_d^2},
\end{equation}
where we use units with $\hbar = m = 1$.

\subsection{Shifted Gaussian Identity}

The key identity for computing kinetic matrix elements is that integrals of polynomials times Gaussians reduce to moments of a Gaussian distribution.

\begin{tcolorbox}[colback=yellow!5,colframe=yellow!50!black,title=Shifted Gaussian Identity]
For a Gaussian overlap $\langle \phi_i | \phi_j \rangle$ with combined width $\Gamma = \alpha_i^* + \alpha_j$, define the effective center:
\begin{equation}
\bar{x} = \frac{\alpha_i^* \tilde{\mu}_i + \alpha_j \tilde{\mu}_j}{\Gamma}
\end{equation}
Then for any polynomial $P(x)$:
\begin{equation}
\langle \phi_i | P(x) | \phi_j \rangle = \langle \phi_i | \phi_j \rangle \cdot \langle P(x) \rangle_{\bar{x}, \sigma^2}
\end{equation}
where $\langle \cdot \rangle_{\bar{x}, \sigma^2}$ denotes expectation with respect to a Gaussian centered at $\bar{x}$ with variance $\sigma^2 = 1/(2\Gamma)$.
\end{tcolorbox}

Specifically, the key moments are:
\begin{align}
\langle 1 \rangle &= 1 \\
\langle x - \mu_j \rangle &= \bar{x} - \mu_j = -\frac{\alpha_i^*}{\Gamma}(\tilde{\mu}_i - \tilde{\mu}_j) = -\frac{\alpha_i^*}{\Gamma} \tilde{y} \\
\langle (x - \mu_j)^2 \rangle &= \frac{1}{2\Gamma} + \left(\frac{\alpha_i^*}{\Gamma}\right)^2 \tilde{y}^2
\end{align}

\subsection{Action of $-\half\partial_x^2$ on a Gaussian}

For a single dimension, the second derivative of $\phi_j(x) = e^{-\alpha_j(x-\mu_j)^2 + ip_j x}$ is:
\begin{align}
\partial_x \phi_j &= \left[ -2\alpha_j(x - \mu_j) + ip_j \right] \phi_j \\
\partial_x^2 \phi_j &= \left[ -2\alpha_j + 4\alpha_j^2(x-\mu_j)^2 - 4i\alpha_j p_j(x-\mu_j) - p_j^2 \right] \phi_j
\end{align}

Therefore:
\begin{equation}
-\half \partial_x^2 \phi_j = \left[ \alpha_j - 2\alpha_j^2(x-\mu_j)^2 + 2i\alpha_j p_j(x-\mu_j) + \half p_j^2 \right] \phi_j
\end{equation}

\subsection{Matrix Element Calculation}

Taking $\langle \phi_i | (-\half\partial_x^2) | \phi_j \rangle$ and using the shifted Gaussian identity:
\begin{align}
T^{(d)}_{ij} &= S^{(d)}_{ij} \left[ \alpha_j - 2\alpha_j^2 \langle(x-\mu_j)^2\rangle + 2i\alpha_j p_j \langle x-\mu_j\rangle + \half p_j^2 \right]
\end{align}

Substituting the moments and simplifying (using $\tilde{k} = \alpha_i^* \alpha_j / \Gamma$ and the definition of $\tilde{y}$), all terms involving $p_j$ cancel with terms from the exponent, yielding:

\begin{tcolorbox}[colback=green!5,colframe=green!50!black,title=Kinetic Matrix $T$]
\begin{equation}
T_{ij} = S_{ij} \cdot \sum_{d=1}^{D} \left[ \tilde{k}_{ij}^{(d)} - 2 (\tilde{k}_{ij}^{(d)})^2 (\tilde{y}_{ij}^{(d)})^2 \right]
\end{equation}
\end{tcolorbox}

Define the per-dimension kinetic factor:
\begin{equation}
F_{ij}^{(d)} = \tilde{k}_{ij}^{(d)} - 2 (\tilde{k}_{ij}^{(d)})^2 (\tilde{y}_{ij}^{(d)})^2.
\end{equation}

Then $T_{ij} = S_{ij} \sum_d F_{ij}^{(d)}$.

\begin{remark}
The simplification works because $\tilde{y}$ already encodes the momentum shifts through $\tilde{\mu}_i$ and $\tilde{\mu}_j$. The final formula depends only on $\tilde{k}$ and $\tilde{y}$, which are the natural "reduced" variables for the Gaussian pair.
\end{remark}

%%%%%%%%%%%%%%%%%%%%%%%%%%%%%%%%%%%%%%%%%%%%%%%%%%%%%%%%%%%%%%%%%%%%%%%%%%%%%%%
\section{Kinetic Squared Matrix $T^2$}
%%%%%%%%%%%%%%%%%%%%%%%%%%%%%%%%%%%%%%%%%%%%%%%%%%%%%%%%%%%%%%%%%%%%%%%%%%%%%%%

The squared kinetic operator is:
\begin{equation}
\hat{T}^2 = \frac{1}{4} \left( \sum_{d=1}^{D} \partial_d^2 \right)^2 = \frac{1}{4} \sum_{d,d'} \partial_d^2 \partial_{d'}^2.
\end{equation}

This splits into diagonal terms ($d = d'$) and off-diagonal terms ($d \neq d'$).

\subsection{Fourth Derivative in One Dimension}

For the diagonal piece, we need $\langle \phi_i | \partial_x^4 | \phi_j \rangle$. Applying $\partial_x^2$ twice to $\phi_j$:
\begin{align}
\partial_x^4 \phi_j &= \Big[ 12\alpha_j^2 - 48\alpha_j^3(x-\mu_j)^2 + 16\alpha_j^4(x-\mu_j)^4 \\
&\quad + (\text{terms with } p_j) \Big] \phi_j \nonumber
\end{align}

Taking the matrix element and using the shifted Gaussian identity for moments up to order 4:
\begin{align}
\langle (x-\mu_j)^4 \rangle &= 3\sigma^4 + 6\sigma^2(\bar{x}-\mu_j)^2 + (\bar{x}-\mu_j)^4
\end{align}
where $\sigma^2 = 1/(2\Gamma)$.

After substitution and simplification (momentum terms cancel as before):

\begin{definition}[Single-Dimension $T^2$ Factor]
\begin{equation}
G_{ij}^{(d)} = 4 (\tilde{k}_{ij}^{(d)})^4 (\tilde{y}_{ij}^{(d)})^4 - 12 (\tilde{k}_{ij}^{(d)})^3 (\tilde{y}_{ij}^{(d)})^2 + 3 (\tilde{k}_{ij}^{(d)})^2
\end{equation}
This corresponds to $\frac{1}{4}\langle \phi_i | \partial_d^4 | \phi_j \rangle / S_{ij}$.
\end{definition}

\subsection{Cross-Dimension Contribution}

For $d \neq d'$, the operators act on different coordinates and the wavefunction factorizes:
\begin{equation}
\langle g_i | \partial_d^2 \partial_{d'}^2 | g_j \rangle = \langle \phi_i^{(d)} | \partial_d^2 | \phi_j^{(d)} \rangle \cdot \langle \phi_i^{(d')} | \partial_{d'}^2 | \phi_j^{(d')} \rangle \cdot \prod_{d'' \neq d, d'} S^{(d'')}_{ij}
\end{equation}

This gives $S_{ij} \cdot F^{(d)} \cdot F^{(d')}$ where $F^{(d)}$ is the kinetic factor from before.

\subsection{Combining All Terms}

The sum over all pairs $(d, d')$:
\begin{equation}
\sum_{d,d'} F^{(d)} F^{(d')} = \left(\sum_d F^{(d)}\right)^2
\end{equation}

But the diagonal terms $d = d'$ contribute $G^{(d)}$, not $(F^{(d)})^2$, so we must correct:
\begin{equation}
\sum_{d \neq d'} F^{(d)} F^{(d')} = \left(\sum_d F^{(d)}\right)^2 - \sum_d (F^{(d)})^2
\end{equation}

\begin{tcolorbox}[colback=green!5,colframe=green!50!black,title=Kinetic Squared Matrix $T^2$]
\begin{equation}
T^2_{ij} = S_{ij} \cdot \left[ \sum_{d=1}^{D} G_{ij}^{(d)} + \left(\sum_{d=1}^{D} F_{ij}^{(d)}\right)^2 - \sum_{d=1}^{D} (F_{ij}^{(d)})^2 \right]
\end{equation}
\end{tcolorbox}

\begin{remark}
The three terms correspond to:
\begin{enumerate}
    \item $\sum_d G^{(d)}$: diagonal fourth-derivative contributions
    \item $(\sum_d F^{(d)})^2$: all pairs of second derivatives
    \item $-\sum_d (F^{(d)})^2$: subtract overcounted diagonal pairs (replaced by $G$)
\end{enumerate}
\end{remark}

%%%%%%%%%%%%%%%%%%%%%%%%%%%%%%%%%%%%%%%%%%%%%%%%%%%%%%%%%%%%%%%%%%%%%%%%%%%%%%%
\section{Polynomial Moments}
%%%%%%%%%%%%%%%%%%%%%%%%%%%%%%%%%%%%%%%%%%%%%%%%%%%%%%%%%%%%%%%%%%%%%%%%%%%%%%%

For polynomial potentials, we need moments of the form $\langle g_i | x_d^k | g_j \rangle$.

\subsection{Generating Moments}

\begin{definition}[Moment Center and Variance]
For a pair $(i,j)$ in dimension $d$:
\begin{align}
P_{ij}^{(d)} &= \frac{\alpha_i^{(d)*} \tilde{\mu}_i^{(d)} + \alpha_j^{(d)} \tilde{\mu}_j^{(d)}}{\Gamma_{ij}^{(d)}} \quad \text{(effective center)} \\[0.3em]
R_{ij}^{(d)} &= \frac{1}{2\Gamma_{ij}^{(d)}} \quad \text{(effective variance)}
\end{align}
\end{definition}

\begin{tcolorbox}[colback=yellow!5,colframe=yellow!50!black,title=Moment Recurrence]
The normalized moments $M_k^{(d)} = \langle g_i | x_d^k | g_j \rangle / S_{ij}$ satisfy:
\begin{align}
M_0^{(d)} &= 1 \\
M_1^{(d)} &= P_{ij}^{(d)} \\
M_{k+1}^{(d)} &= P_{ij}^{(d)} M_k^{(d)} + k \, R_{ij}^{(d)} M_{k-1}^{(d)}
\end{align}
\end{tcolorbox}

This is the Hermite polynomial recurrence for moments of a Gaussian distribution.

\subsection{Multi-Dimensional Polynomial Expectation}

For a monomial $\bm{x}^{\bm{n}} = \prod_{d=1}^{D} x_d^{n_d}$:
\begin{equation}
\langle g_i | \bm{x}^{\bm{n}} | g_j \rangle = S_{ij} \prod_{d=1}^{D} M_{n_d}^{(d)}.
\end{equation}

%%%%%%%%%%%%%%%%%%%%%%%%%%%%%%%%%%%%%%%%%%%%%%%%%%%%%%%%%%%%%%%%%%%%%%%%%%%%%%%
\section{Polynomial Potential $V_{\text{poly}}$}
%%%%%%%%%%%%%%%%%%%%%%%%%%%%%%%%%%%%%%%%%%%%%%%%%%%%%%%%%%%%%%%%%%%%%%%%%%%%%%%

A polynomial potential is:
\begin{equation}
V_{\text{poly}}(\bm{x}) = \sum_{\bm{n}} v_{\bm{n}} \, \bm{x}^{\bm{n}},
\end{equation}
where the sum is over a finite set of multi-indices $\bm{n} = (n_1, \ldots, n_D)$.

\begin{tcolorbox}[colback=red!5,colframe=red!50!black,title=Polynomial Potential Matrix]
\begin{equation}
(V_{\text{poly}})_{ij} = \sum_{\bm{n}} v_{\bm{n}} \, S_{ij} \prod_{d=1}^{D} M_{n_d}^{(d)}
\end{equation}
\end{tcolorbox}

%%%%%%%%%%%%%%%%%%%%%%%%%%%%%%%%%%%%%%%%%%%%%%%%%%%%%%%%%%%%%%%%%%%%%%%%%%%%%%%
\section{Gaussian Potential $V_{\text{gauss}}$}
%%%%%%%%%%%%%%%%%%%%%%%%%%%%%%%%%%%%%%%%%%%%%%%%%%%%%%%%%%%%%%%%%%%%%%%%%%%%%%%

The Gaussian potential is a sum of Gaussian terms:
\begin{equation}
V_{\text{gauss}}(\bm{x}) = \sum_{k=1}^{m} c_k \prod_{d=1}^{D} \exp\left[ -\alpha_k^{(d)} (x_d - \mu_k^{(d)})^2 \right],
\end{equation}
where $c_k$ are coefficients and $\alpha_k^{(d)} > 0$ (real, purely Gaussian width).

\subsection{Convolution with Basis Functions}

The matrix element $\langle g_i | V_k | g_j \rangle$ for a single Gaussian potential term involves combining three Gaussians. The key insight is that the product of two Gaussians is again a Gaussian.

\begin{definition}[Combined Gaussian Parameters]
For basis pair $(i,j)$ and potential $k$:
\begin{align}
A_{ij}^{(d)} &= \alpha_i^{(d)*} + \alpha_j^{(d)} \quad \text{(from basis overlap)} \\
B_{ij}^{(d)} &= 2(\alpha_i^{(d)*} \mu_i^{(d)} + \alpha_j^{(d)} \mu_j^{(d)}) + i(p_j^{(d)} - p_i^{(d)}) \\
C_{ij}^{(d)} &= -\alpha_i^{(d)*} (\mu_i^{(d)})^2 - \alpha_j^{(d)} (\mu_j^{(d)})^2 + i(\mu_i^{(d)} p_i^{(d)} - \mu_j^{(d)} p_j^{(d)})
\end{align}
When adding potential $k$:
\begin{align}
\tilde{A}^{(d)} &= A_{ij}^{(d)} + \alpha_k^{(d)} \\
\tilde{B}^{(d)} &= B_{ij}^{(d)} + 2\alpha_k^{(d)} \mu_k^{(d)} \\
\tilde{C}^{(d)} &= C_{ij}^{(d)} - \alpha_k^{(d)} (\mu_k^{(d)})^2
\end{align}
\end{definition}

\subsection{Gaussian Potential Matrix Element}

Completing the square:
\begin{equation}
\text{center} = \frac{\tilde{B}^{(d)}}{2\tilde{A}^{(d)}}, \quad
\text{constant} = \tilde{C}^{(d)} + \frac{(\tilde{B}^{(d)})^2}{4\tilde{A}^{(d)}}.
\end{equation}

\begin{tcolorbox}[colback=red!5,colframe=red!50!black,title=Gaussian Potential Matrix]
\begin{equation}
(V_{\text{gauss}})_{ij} = \sum_{k=1}^{m} c_k \exp\left[ \sum_{d=1}^{D} \left( \text{constant}_k^{(d)} - \half \ln \tilde{A}_k^{(d)} - \half E_{\text{sub}}^{(d)} \right) \right]
\end{equation}
where $E_{\text{sub}}^{(d)}$ is the normalization subtraction term (diagonal elements of basis overlap).
\end{tcolorbox}

%%%%%%%%%%%%%%%%%%%%%%%%%%%%%%%%%%%%%%%%%%%%%%%%%%%%%%%%%%%%%%%%%%%%%%%%%%%%%%%
\section{Potential Squared $V^2$}
%%%%%%%%%%%%%%%%%%%%%%%%%%%%%%%%%%%%%%%%%%%%%%%%%%%%%%%%%%%%%%%%%%%%%%%%%%%%%%%

The squared potential involves:
\begin{equation}
\langle g_i | V^2 | g_j \rangle = \langle g_i | V_{\text{poly}}^2 + V_{\text{gauss}}^2 + 2 V_{\text{poly}} V_{\text{gauss}} | g_j \rangle.
\end{equation}

\subsection{Polynomial Squared}

For $V_{\text{poly}}^2$, we simply compute the squared polynomial by multiplying out all terms and collecting coefficients for each monomial $\bm{x}^{\bm{m}}$.

If $V_{\text{poly}} = \sum_{\bm{n}} v_{\bm{n}} \bm{x}^{\bm{n}}$, then:
\begin{equation}
V_{\text{poly}}^2 = \sum_{\bm{n}, \bm{n}'} v_{\bm{n}} v_{\bm{n}'} \bm{x}^{\bm{n} + \bm{n}'}.
\end{equation}

\subsection{Gaussian Squared}

For two Gaussian potential terms $k$ and $k'$, their product is:
\begin{equation}
V_k V_{k'} = c_k c_{k'} \prod_{d=1}^D \exp\left[ -(\alpha_k^{(d)} + \alpha_{k'}^{(d)})(x_d - \mu_{kk'}^{(d)})^2 + \text{const} \right],
\end{equation}
which is again a Gaussian. The squared Gaussian potential is precomputed by forming all $m^2$ product terms.

\subsection{Polynomial-Gaussian Cross Terms}

For $V_{\text{poly}} V_{\text{gauss}}$, each Gaussian potential modifies the effective center and variance for the moment calculation:
\begin{align}
P_k^{(d)} &= \frac{\tilde{B}_k^{(d)}}{2\tilde{A}_k^{(d)}} \quad \text{(new center)} \\
R_k^{(d)} &= \frac{1}{2\tilde{A}_k^{(d)}} \quad \text{(new variance)}
\end{align}

The moments are recomputed with these modified parameters, then contracted with the polynomial coefficients.

%%%%%%%%%%%%%%%%%%%%%%%%%%%%%%%%%%%%%%%%%%%%%%%%%%%%%%%%%%%%%%%%%%%%%%%%%%%%%%%
\section{Kinetic-Potential Cross Terms $TV + VT$}
%%%%%%%%%%%%%%%%%%%%%%%%%%%%%%%%%%%%%%%%%%%%%%%%%%%%%%%%%%%%%%%%%%%%%%%%%%%%%%%

The Hamiltonian squared contains:
\begin{equation}
H^2 = T^2 + V^2 + TV + VT.
\end{equation}

Since $T$ and $V$ generally don't commute, we need both $TV$ and $VT$.

\subsection{Hermiticity Relation}

Since both $T$ and $V$ are Hermitian:
\begin{equation}
(VT)_{ij}^* = \langle g_j | VT | g_i \rangle^* = \langle g_i | TV | g_j \rangle = (TV)_{ij}.
\end{equation}

Thus: $\boxed{TV = (VT)^\dagger}$.

We only need to compute $VT$ and obtain $TV$ by conjugate transpose.

\subsection{Polynomial VT}

Acting with $T = -\half \nabla^2$ on $|g_j\rangle$ and taking the matrix element with $V_{\text{poly}}$:

\begin{tcolorbox}[colback=purple!5,colframe=purple!50!black,title=Polynomial $VT$ Matrix]
For monomial $\bm{x}^{\bm{n}}$ with coefficient $v_{\bm{n}}$, the contribution to $(VT)_{ij}$ involves:
\begin{equation}
\text{base}^{(d)} = -2(\tilde{k}^{(d)})^2 (\tilde{y}^{(d)})^2 + \tilde{k}^{(d)}
\end{equation}
and correction terms from the polynomial:
\begin{equation}
\text{vt\_coeff}^{(d)} = \text{base}^{(d)} - 2 n_d \tilde{k}^{(d)} \tilde{y}^{(d)} \beta_j^{(d)} \frac{M_{n_d-1}}{M_{n_d}} - \half n_d(n_d-1) (\beta_j^{(d)})^2 \frac{M_{n_d-2}}{M_{n_d}}
\end{equation}
where $\beta_j^{(d)} = \alpha_j^{(d)} / \Gamma_{ij}^{(d)}$.

Then:
\begin{equation}
(VT)_{ij} = \sum_{\bm{n}} v_{\bm{n}} \, S_{ij} \prod_{d=1}^D M_{n_d}^{(d)} \cdot \sum_{d=1}^D \text{vt\_coeff}^{(d)}
\end{equation}
\end{tcolorbox}

\subsection{Gaussian VT}

For Gaussian potentials, we compute expectation values of the Laplacian in the presence of the Gaussian weight.

\begin{definition}[Gaussian-Weighted Moments]
With potential $k$ present:
\begin{align}
\langle (x - \mu_k)^2 \rangle_k &= R_k + (P_k - \mu_k)^2 \\
\langle (x - \mu_j)(x - \mu_k) \rangle_k &= R_k + (P_k - \mu_j)(P_k - \mu_k)
\end{align}
\end{definition}

\begin{tcolorbox}[colback=purple!5,colframe=purple!50!black,title=Gaussian $VT$ Matrix]
For each potential term $k$:
\begin{align}
\text{term}_1 &= -\half \sum_d \left[ 4\alpha_k^2 \langle(x-\mu_k)^2\rangle_k - 2\alpha_k \right] \\
\text{term}_2 &= -\sum_d (-2\alpha_k)\left[ (-2\alpha_j) \text{Cov}_{kj}^{(d)} + i p_j \langle x-\mu_k \rangle_k \right] \\
\text{term}_3 &= -\sum_d \left[ -2\alpha_j + 4\alpha_j^2 \langle(x-\mu_j)^2\rangle_k - 4i\alpha_j p_j \langle x-\mu_j\rangle_k - p_j^2 \right]
\end{align}
where $\text{Cov}_{kj}^{(d)} = R_k + (P_k - \mu_k)(P_k - \mu_j)$.

Then:
\begin{equation}
(VT)_{ij} = \sum_{k=1}^{m} c_k \cdot \text{prefactor}_k \cdot (\text{term}_1 + \text{term}_2 + \text{term}_3)
\end{equation}
\end{tcolorbox}

The full cross term is:
\begin{equation}
TV + VT = VT + VT^\dagger.
\end{equation}

%%%%%%%%%%%%%%%%%%%%%%%%%%%%%%%%%%%%%%%%%%%%%%%%%%%%%%%%%%%%%%%%%%%%%%%%%%%%%%%
\section{The Rothe Method}
%%%%%%%%%%%%%%%%%%%%%%%%%%%%%%%%%%%%%%%%%%%%%%%%%%%%%%%%%%%%%%%%%%%%%%%%%%%%%%%

\subsection{Time-Stepping Formulation}

The Rothe method approximates the time evolution by minimizing the residual of the implicit Euler scheme:
\begin{equation}
\left\| \frac{\Psi^{n+1} - \Psi^n}{\Delta t} + i H \Psi^{n+1} \right\|^2 = \min.
\end{equation}

\subsection{Matrix Form}

Expanding in the Gaussian basis $\Psi^{n+1} = \sum_j c_j^{n+1} g_j^{n+1}$:

\begin{tcolorbox}[colback=gray!5,colframe=gray!50!black,title=Rothe Error Functional]
\begin{equation}
\mathcal{E}[\bm{\theta}^{n+1}] = \bm{c}^{n+1\dagger} \left( S + i\Delta t H + \frac{\Delta t^2}{2} H^2 \right) \bm{c}^{n+1} - 2\re\left( \bm{c}^{n+1\dagger} S_{01} \bm{c}^n \right) + \langle \Psi^n | \Psi^n \rangle
\end{equation}
where:
\begin{itemize}
    \item $S$ is the overlap matrix of new basis functions
    \item $H$ is the Hamiltonian matrix in the new basis
    \item $H^2$ is the squared Hamiltonian matrix
    \item $S_{01}$ is the overlap between old and new bases
\end{itemize}
\end{tcolorbox}

\subsection{Optimal Coefficients}

For fixed Gaussian parameters $\bm{\theta}^{n+1}$, the optimal coefficients are:
\begin{equation}
\bm{c}^{n+1} = \left( S + i\Delta t H + \frac{\Delta t^2}{2} H^2 \right)^{-1} S_{01} \bm{c}^n.
\end{equation}

The optimization then proceeds over the Gaussian parameters $\bm{\theta}^{n+1}$.

%%%%%%%%%%%%%%%%%%%%%%%%%%%%%%%%%%%%%%%%%%%%%%%%%%%%%%%%%%%%%%%%%%%%%%%%%%%%%%%
\section{Summary of Intermediates}
%%%%%%%%%%%%%%%%%%%%%%%%%%%%%%%%%%%%%%%%%%%%%%%%%%%%%%%%%%%%%%%%%%%%%%%%%%%%%%%

\begin{table}[h]
\centering
\begin{tabular}{|c|l|c|}
\hline
\textbf{Symbol} & \textbf{Description} & \textbf{Shape} \\
\hline
$\alpha_i^{(d)}$ & Complex width: $a_i^2 + i b_i$ & $(n, D)$ \\
$\Gamma_{ij}^{(d)}$ & Sum: $\alpha_i^* + \alpha_j$ & $(n, n, D)$ \\
$\tilde{y}_{ij}^{(d)}$ & Shifted center difference & $(n, n, D)$ \\
$\tilde{k}_{ij}^{(d)}$ & $\alpha_i^* \alpha_j / \Gamma_{ij}$ & $(n, n, D)$ \\
$P_{ij}^{(d)}$ & Effective center for moments & $(n, n, D)$ \\
$R_{ij}^{(d)}$ & Effective variance: $1/(2\Gamma)$ & $(n, n, D)$ \\
$\beta_j^{(d)}$ & Ratio: $\alpha_j / \Gamma_{ij}$ & $(n, n, D)$ \\
$M_k^{(d)}$ & $k$-th moment (normalized) & $(n, n, D, K_{\max}+1)$ \\
\hline
\end{tabular}
\caption{Key intermediate quantities computed once per basis update.}
\end{table}

%%%%%%%%%%%%%%%%%%%%%%%%%%%%%%%%%%%%%%%%%%%%%%%%%%%%%%%%%%%%%%%%%%%%%%%%%%%%%%%
\section{Implementation Notes}
%%%%%%%%%%%%%%%%%%%%%%%%%%%%%%%%%%%%%%%%%%%%%%%%%%%%%%%%%%%%%%%%%%%%%%%%%%%%%%%

\subsection{Memory Efficiency}

\begin{itemize}
    \item Gaussian potential sums are computed via \texttt{lax.scan} to avoid materializing $(m, n, n, D)$ tensors.
    \item Custom VJP rules provide memory-efficient gradients for the backward pass.
    \item Moments are precomputed up to $K_{\max}$ determined by the polynomial degree.
\end{itemize}

\subsection{GPU Optimization}

\begin{itemize}
    \item The Gaussian kinetic cross terms are vectorized over the $m$ potential terms.
    \item Avoid \texttt{lax.fori\_loop} for small $m$ as it serializes GPU work.
    \item Use \texttt{tensordot}-style contractions rather than explicit loops.
\end{itemize}

\subsection{Numerical Stability}

\begin{itemize}
    \item All computations use \texttt{complex128} for accuracy.
    \item Overlap matrix is computed in log-space to avoid underflow.
    \item Safe division with $\epsilon$ guards for moment ratios.
\end{itemize}

\end{document}
